\documentclass[11pt,a4paper,UTF8]{ctexart}
\usepackage{geometry}
\geometry{left=18mm,right=18mm,top=24mm,bottom=24mm,headheight=22pt,footskip=12mm}
\usepackage{amsmath,amssymb,mathtools,bm,esint,extarrows,centernot}
\usepackage{xcolor}
\usepackage{tcolorbox}
\tcbuselibrary{skins,breakable}
\usepackage{fancyhdr}
\usepackage{titlesec}
\usepackage{hyperref}
\usepackage{tikz}
\usepackage{booktabs}
\usepackage{tabularx}

% --- Color Palette (Purple & DeepPink Academic Theme) ---
\definecolor{DeepPurple}{HTML}{4C1D95}
\definecolor{TitlePurple}{HTML}{581C87}
\definecolor{SubPurple}{HTML}{6B21A8}
\definecolor{DeepPink}{HTML}{FF1493}
\definecolor{LightPinkBg}{HTML}{FFF5F8}
\definecolor{LightPurpleBg}{HTML}{F8F5FF}
\definecolor{GoldAccent}{HTML}{D97706}
\definecolor{DarkText}{HTML}{1F2937}

\hypersetup{
    colorlinks=true,
    linkcolor=TitlePurple,
    citecolor=DeepPink,
    urlcolor=SubPurple,
    pdfborder={0 0 0}
}

% --- Typography & Spacing ---
\linespread{1.3}
\setlength{\parskip}{0.35em plus 0.1em minus 0.1em}
\setlength{\parindent}{0pt}

% --- Section Titles Styling ---
\titleformat{\section}
  {\Large\bfseries\color{TitlePurple}}{\thesection}{1em}{}
  [\vspace{1mm}{\color{TitlePurple!30}\hrule height 1pt}]
\titleformat{\subsection}
  {\large\bfseries\color{SubPurple}}{\thesubsection}{0.8em}{}
\titleformat{\subsubsection}
  {\normalsize\bfseries\color{DeepPurple}}{\thesubsubsection}{0.6em}{}

% --- tcolorbox Environments ---
% 1. Theorem / Formula / Method / Analysis Box (Deep Purple)
\newtcolorbox{thmbox}[1][]{
    enhanced,
    breakable,
    colback=LightPurpleBg,
    colframe=TitlePurple,
    arc=3pt,
    boxrule=0pt,
    leftrule=3.5pt,
    left=10pt,
    right=10pt,
    top=8pt,
    bottom=8pt,
    title=#1,
    coltitle=white,
    colbacktitle=TitlePurple,
    fonttitle=\bfseries\small,
    attach boxed title to top left={yshift=-2mm, xshift=4mm},
    boxed title style={boxrule=0pt, arc=2pt}
}

% 2. Solution / Formula / Derivation Box (DeepPink)
\newtcolorbox{solbox}[1][]{
    enhanced,
    breakable,
    colback=LightPinkBg,
    colframe=DeepPink,
    arc=3pt,
    boxrule=0pt,
    leftrule=3.5pt,
    left=10pt,
    right=10pt,
    top=8pt,
    bottom=8pt,
    title=#1,
    coltitle=white,
    colbacktitle=DeepPink,
    fonttitle=\bfseries\small,
    attach boxed title to top left={yshift=-2mm, xshift=4mm},
    boxed title style={boxrule=0pt, arc=2pt}
}

% 3. Learning Goals / Summary Box
\newtcolorbox{targetbox}[1][]{
    enhanced,
    breakable,
    colback=LightPurpleBg,
    colframe=DeepPurple,
    arc=3pt,
    boxrule=1pt,
    left=10pt,
    right=10pt,
    top=8pt,
    bottom=8pt,
    title=#1,
    coltitle=white,
    colbacktitle=DeepPurple,
    fonttitle=\bfseries\small,
    attach boxed title to top left={yshift=-2mm, xshift=4mm},
    boxed title style={boxrule=0pt, arc=2pt}
}

\begin{document}

% ------------------- 讲义封面 (Cover Page) -------------------
\begin{titlepage}
\thispagestyle{empty}
\begin{tikzpicture}[remember picture,overlay]
    \fill[DeepPurple] (current page.north west) rectangle ([yshift=-7cm]current page.north east);
    \draw[DeepPink, line width=3.5pt] ([yshift=-7cm]current page.north west) -- ([yshift=-7cm]current page.north east);
    \draw[GoldAccent, line width=1pt] ([yshift=-7.12cm]current page.north west) -- ([yshift=-7.12cm]current page.north east);
    
    \node[anchor=north east, opacity=0.12, text=white] at ([xshift=-1.5cm, yshift=-0.5cm]current page.north east) {
        \fontsize{70}{70}\selectfont\bfseries 2027
    };
\end{tikzpicture}

\vspace*{0.8cm}
\begin{center}
    {\color{white}\fontsize{26}{32}\selectfont\bfseries 2027 考研数学(二) 核心讲义}\\[0.6cm]
    {\color{white}\fontsize{13}{18}\selectfont\bfseries 全国硕士研究生招生考试 · 统考数学(二) 核心精讲全书}\\[1.5cm]
\end{center}

\vspace{3.5cm}
\begin{center}
    {\color{GoldAccent}\large\bfseries $\blacktriangleright$ 基础强化 · 核心题解 · 考点深水区 $\blacktriangleleft$}\\[0.8cm]
    {\color{TitlePurple}\fontsize{24}{28}\selectfont\bfseries 第四章 一元函数微分学的计算}\\[0.6cm]
    {\color{DeepPink}\large\bfseries —— 考研高阶冲刺与深水区专攻 ——}\\[1.5cm]
\end{center}

\vfill

\begin{center}
\begin{tcolorbox}[
    colback=white,
    colframe=DeepPurple,
    width=0.88\textwidth,
    arc=4pt,
    boxrule=1.2pt,
    halign=center,
    drop shadow=gray!30
]
    \vspace{3mm}
    {\bfseries\large 编著团队：EscoffierZhou}\\[2.5mm]
    {\bfseries\color{TitlePurple} 目标院校：中国科学院大学 (UCAS) 计算机专硕 (22408)}\\[2.5mm]
    {\small\color{gray} 备考铁律：手算真题数字 · 错题白纸重做 · 408 客观题为王}\\[2.5mm]
    {\small 编制版本：2027 届首发版 · \today}
    \vspace{2mm}
\end{tcolorbox}
\end{center}
\vspace{0.8cm}
\end{titlepage}

% ------------------- 目录与页面主体配置 -------------------
\pagestyle{fancy}
\fancyhf{}
\fancyhead[L]{\small\color{gray}2027 考研数学(二) · 核心讲义系列}
\fancyhead[R]{\small\color{TitlePurple}\nouppercase{\leftmark}}
\fancyfoot[C]{\small\color{gray}— 第 \thepage\ 页 —}
\fancyfoot[R]{\small\color{gray}UCAS 22408}
\renewcommand{\headrulewidth}{0.5pt}

\pagenumbering{Roman}
\tableofcontents
\newpage
\pagenumbering{arabic}


\subsubsection{1. 变上限积分函数的求导（数二绝对的核心大题考点）}

数二不考三重积分和曲面积分，因此变限积分的求导与综合题是数二拉开分差的命门，必须补入以下 SOP：

-   \textbf{核心公式（复合上限）}：
\begin{equation*}
\frac{\mathrm{d}}{\mathrm{d}x} \int_{\alpha(x)}^{\beta(x)} f(t)\,\mathrm{d}t = f[\beta(x)]\beta'(x) - f[\alpha(x)]\alpha'(x)
\end{equation*}

-   \textit{}被积函数含自变量 $x$ 的“剥离/拆分”求导法\textit{}： 遇到 $F(x) = \int_0^x f(x, t)\,\mathrm{d}t$，\textit{}绝不能直接把 $x$ 往里套\textit{}：
    -   \textbf{类型 A（线性可拆）}：$\int_0^x (x - t) f(t)\,\mathrm{d}t = x\int_0^x f(t)\,\mathrm{d}t - \int_0^x t f(t)\,\mathrm{d}t$，拆开后再用乘积求导法则。
    -   \textbf{类型 B（复合不可拆，必须换元）}：如 $\int_0^x f(x - t)\,\mathrm{d}t$ 或 $\int_0^x t f(x^2 - t^2)\,\mathrm{d}t$。必须先令 $u = x - t$ 将参变量 $x$ 换元到积分限上去，消去被积函数内部的 $x$，再进行求导。


\subsubsection{2. 极坐标方程下的求导与切线（数二专项）}

数二每隔几年必考极坐标曲线 $r = r(\theta)$ 的几何应用与求导计算：

-   \textbf{转化为参数方程}：$\begin{cases} x = r(\theta)\cos\theta \\ y = r(\theta)\sin\theta \end{cases}$

-   \textbf{切线斜率}：


\begin{equation*}
\frac{\mathrm{d}y}{\mathrm{d}x} = \frac{\frac{\mathrm{d}y}{\mathrm{d}\theta}}{\frac{\mathrm{d}x}{\mathrm{d}\theta}} = \frac{r'(\theta)\sin\theta + r(\theta)\cos\theta}{r'(\theta)\cos\theta - r(\theta)\sin\theta}
\end{equation*}


\subsubsection{3. 隐函数存在定理与多元偏导交叉法（降维打击工具）}

-   对于方程 $F(x, y) = 0$，强化阶段除了两边求导，必须熟练掌握偏导数商公式：


\begin{equation*}
\frac{\mathrm{d}y}{\mathrm{d}x} = -\frac{F'_x}{F'_y} \quad (F'_y \neq 0)
\end{equation*}

-   \textbf{二阶隐函数导数的代入提速策略}：不要推导冗长的显式二阶公式，牢记对等式 $F'_x + F'_y \cdot y' = 0$ 两端继续对 $x$ 求导，代入点值直接解线性方程。


\subsubsection{4. 抽象复合函数的高阶链式展开}

强化题极少直接给初等函数让你求导，全是多层抽象复合函数：

-   设 $z = f(x^2 - y, \, g(x))$，求 $\frac{\mathrm{d}z}{\mathrm{d}x}$。
-   必须训练规范的\textbf{位置标记法}（如 $f'_1$ 表示对第 1 个位置求偏导），杜绝变量名混淆。


\subsubsection{5. “非零因子”极限与导数定义的嫁接计算}

求极限题与导数定义的深度结合，典型题型：


\begin{equation*}
\lim_{x \to 0} \frac{f(x) - \sin x}{x^2} = 2 \implies \text{反求 } f(0), f'(0), f''(0)
\end{equation*}

-   这是数二最爱考的压轴选填题，必须掌握“利用局部泰勒展开/导数定义逆推原函数导数值”的计算方法。

\end{document}
